\section{Day 8: More Distributions (Sep 29, 2025)}

\begin{definition}[Geometric Distribution]
    $P(Z = k) = (1 - \theta)^k \theta$ 
\end{definition}

Intuitively, this can be thought of as $k$ failed attempts at something with $\theta$ probability each time, with success on the $k + 1$-th attempt. Some textbooks count the number of attempts to be up to and including the first success, not here though. 

\begin{simplethm}
    If $X \sim \text{Geometric}(\theta)$ for any $0 < \theta < 1$, with $m \in \mathbb{N}$, then $P(X \geq m) = (1 - \theta)^m > 0$, yet $P(X = \infty) = 0$.
\end{simplethm}

\begin{simplethm}[Memoryless Property]
    If $X \sim \text{Geometric}(\theta)$, $a, b \in \mathbb{N}$ then 
\[
    P(X \geq a + b \mid X \geq a) = P(X \geq b)
\]
\end{simplethm}

\noindent An analogy is that the probability of you having to wait an additional $b$ minutes, is `independent' of you having already waited $a$ minutes. 

\begin{definition}[Poisson Distribution]
    Given $\lambda$, and $k \in \mathbb{N} \cup \{ 0 \}$,
    \[
    P(k) = e^{-\lambda} \frac{\lambda^k}{k!}
    \]
\end{definition}

\noindent Suppose we have $n$ houses in a city, with each house independently having a $\theta$ probability of a fire occurring. The number of houses on fire $\lambda = n \theta$. We can compute the probability of having $k$ houses on fire as follows:
\[
P(\text{\# fires} = k) = \binom{n}{k} (1-\theta)^{n-k} \theta^{k}
\]

\noindent We will be skipping the Negative-Binomial($r, \theta$) and Hypergeometric($N$, $M$, $n$) distributions, and cover them next year.

\begin{theorem}
If $X$ is a discrete random variable, with values $x_1, x_2, \cdots$ and corresponding probabilities $p_1, p_2, \cdots$ and $B$ is any event, then
\[
    P(B) = \sum_i P(X = x_i) P(B \mid X = x_i) = \sum_i p_i P(B \mid X = x_i)
\]
\end{theorem}

\noindent This follows from \ref{thm:ltpcond}, since $X = x_i$ for some $i$ will always occur, by definition of $x_i$, thus they union to $S$. For pairwise disjointedness $X$ can only take on one value at a time, hence this is a valid partition.

\begin{remark}
For now we say that $\binom{0}{3}$ is undefined. There is a way to define it using negative binomials, but we're not interested in that just yet.
\end{remark}

\subsection{Continuous Random Variables}

\begin{definition}[Continuous]
$X$ is continuous if for all $x \in S$, $P(X = x) = 0$.
\end{definition}

\begin{definition}[Density Function]
    A density function is `any' $f: \mathbb{R} \to \mathbb{R}$ with $f(x) \geq 0$, and $\int_{-\infty}^{\infty} f(x) \, dx  = 1$.
\end{definition}

Given any density function $f$, we can define $X$, where $P(a \leq X \leq b) = \int_a^b f(x) \, dx$ for $a \leq b$. Immediately $X$ is continuous, since for any $x \in S$, have $P(X = x) = \int_x^x f(t) \, dt = 0$. We call such an $X$, \textbf{absolutely continuous} and write $f_X(x)$ to denote that it's $X$'s density function.
